% GENERATED FILE. Edit canonical lessons, then run npm run generate:books.
% canonical-source-hash: fnv1a64:eb89087113503a1b
% canonical-lessons: SA-C120-mudra, SA-C120-patralayah, SA-C120-lekhah, SA-C120-granthah, SA-C120-citram

\chapter{Coin, Post office, Letter, Book, Picture}
\label{ch:sa-coin-post-office-letter-book-picture}
\hlchaptermodality{120}

\begin{chapteropening}
\textbf{By the end of this chapter:} \emph{I can say a coin, a post office, a letter, a book and a picture.}

Complete the last lesson of chapter 120: i can say a coin, a post office, a letter, a book and a picture.
\end{chapteropening}

\section[\texorpdfstring{\sk{मुद्रा} (mudrā)}{मुद्रा (mudrā)}]{\sk{मुद्रा} (mudrā) — a coin, a seal}

\label{lesson:SA-C120-mudra}



\begin{quote}
Before the new one: say the Sanskrit for a bottle, a flask, then the Sanskrit for a small umbrella.
\end{quote}

\subsection*{You'll want to know: \sk{मुद्रा}}
\textbf{\sk{मुद्रा}} — \emph{mudrā} — \textquotedblleft{}a coin, a seal\textquotedblright{}.

\begin{grammarlens}[title={What you've built}]
One more word for this chapter.
\end{grammarlens}

\subsection*{Your turn}
\begin{itemize}
\raggedright
  \item \emph{Say it:} \emph{mudrā}
  \item \emph{Say it:} \emph{mudrā}, once more
  \item \emph{Say it:} say \emph{mudrā}
  \item \emph{Recall:} say the Sanskrit for a bottle, a flask, then the Sanskrit for a small umbrella, then say \emph{mudrā} again
\end{itemize}

\begin{tcolorbox}[breakable,colback=teal!4,colframe=teal!35!black,title={Before you move on}]
What does \sk{मुद्रा} mean? (\textquotedblleft{}A coin, a seal\textquotedblright{}.) Say it once more.
\end{tcolorbox}
\section[\texorpdfstring{\sk{पत्रालयः} (patrālayaḥ)}{पत्रालयः (patrālayaḥ)}]{\sk{पत्रालयः} (patrālayaḥ) — a post office}

\label{lesson:SA-C120-patralayah}



\begin{quote}
Before the new one: say the Sanskrit for a small umbrella, then the Sanskrit for a coin.
\end{quote}

\subsection*{You'll want to know: \sk{पत्रालयः}}
\textbf{\sk{पत्रालयः}} — \emph{patrālayaḥ} — \textquotedblleft{}a post office\textquotedblright{}.

\begin{grammarlens}[title={What you've built}]
One more word for this chapter.
\end{grammarlens}

\subsection*{Your turn}
\begin{itemize}
\raggedright
  \item \emph{Say it:} \emph{patrālayaḥ}
  \item \emph{Say it:} \emph{patrālayaḥ}, once more
  \item \emph{Say it:} say \emph{patrālayaḥ}
  \item \emph{Recall:} say the Sanskrit for a small umbrella, then the Sanskrit for a coin, then say \emph{patrālayaḥ} again
\end{itemize}

\begin{tcolorbox}[breakable,colback=teal!4,colframe=teal!35!black,title={Before you move on}]
What does \sk{पत्रालयः} mean? (\textquotedblleft{}A post office\textquotedblright{}.) Say it once more.
\end{tcolorbox}
\section[\texorpdfstring{\sk{लेखः} (lekhaḥ)}{लेखः (lekhaḥ)}]{\sk{लेखः} (lekhaḥ) — a letter, an article}

\label{lesson:SA-C120-lekhah}



\begin{quote}
Before the new one: say the Sanskrit for a coin, then the Sanskrit for a post office.
\end{quote}

\subsection*{You'll want to know: \sk{लेखः}}
\textbf{\sk{लेखः}} — \emph{lekhaḥ} — \textquotedblleft{}a letter, an article\textquotedblright{}.

\begin{grammarlens}[title={What you've built}]
One more word for this chapter.
\end{grammarlens}

\subsection*{Your turn}
\begin{itemize}
\raggedright
  \item \emph{Say it:} \emph{lekhaḥ}
  \item \emph{Say it:} \emph{lekhaḥ}, once more
  \item \emph{Say it:} say \emph{lekhaḥ}
  \item \emph{Recall:} say the Sanskrit for a coin, then the Sanskrit for a post office, then say \emph{lekhaḥ} again
\end{itemize}

\begin{tcolorbox}[breakable,colback=teal!4,colframe=teal!35!black,title={Before you move on}]
What does \sk{लेखः} mean? (\textquotedblleft{}A letter, an article\textquotedblright{}.) Say it once more.
\end{tcolorbox}
\section[\texorpdfstring{\sk{ग्रन्थः} (granthaḥ)}{ग्रन्थः (granthaḥ)}]{\sk{ग्रन्थः} (granthaḥ) — a book, a text}

\label{lesson:SA-C120-granthah}



\begin{quote}
Before the new one: say the Sanskrit for a post office, then the Sanskrit for a letter.
\end{quote}

\subsection*{You'll want to know: \sk{ग्रन्थः}}
\textbf{\sk{ग्रन्थः}} — \emph{granthaḥ} — \textquotedblleft{}a book, a text\textquotedblright{}.

\begin{grammarlens}[title={What you've built}]
One more word for this chapter.
\end{grammarlens}

\subsection*{Your turn}
\begin{itemize}
\raggedright
  \item \emph{Say it:} \emph{granthaḥ}
  \item \emph{Say it:} \emph{granthaḥ}, once more
  \item \emph{Say it:} say \emph{granthaḥ}
  \item \emph{Recall:} say the Sanskrit for a post office, then the Sanskrit for a letter, then say \emph{granthaḥ} again
\end{itemize}

\begin{tcolorbox}[breakable,colback=teal!4,colframe=teal!35!black,title={Before you move on}]
What does \sk{ग्रन्थः} mean? (\textquotedblleft{}A book, a text\textquotedblright{}.) Say it once more.
\end{tcolorbox}
\section[\texorpdfstring{\sk{चित्रम्} (citram)}{चित्रम् (citram)}]{\sk{चित्रम्} (citram) — a picture}

\label{lesson:SA-C120-citram}



\begin{quote}
Before the new one: say the Sanskrit for a letter, then the Sanskrit for a book.
\end{quote}

\subsection*{You'll want to know: \sk{चित्रम्}}
\textbf{\sk{चित्रम्}} — \emph{citram} — \textquotedblleft{}a picture\textquotedblright{}.

\begin{grammarlens}[title={What you've built}]
One more word for this chapter.
\end{grammarlens}

\subsection*{Your turn}
\begin{itemize}
\raggedright
  \item \emph{Say it:} \emph{citram}
  \item \emph{Say it:} \emph{citram}, once more
  \item \emph{Say it:} say \emph{citram}
  \item \emph{Recall:} say the Sanskrit for a letter, then the Sanskrit for a book, then say \emph{citram} again
\end{itemize}

\begin{tcolorbox}[breakable,colback=teal!4,colframe=teal!35!black,title={Before you move on}]
What does \sk{चित्रम्} mean? (\textquotedblleft{}A picture\textquotedblright{}.) Say it once more.
\end{tcolorbox}
